Al entrenar un modelo inverso basado en redes neuronales profundas con datos EEG simulados, donde las ubicaciones de las fuentes neuronales son conocidas, es posible mejorar la precisión en la detección del origen de la señal cerebral. Este enfoque supera a métodos clásicos como MNE, sLORETA y VARETA en la estimación espacio-temporal de las fuentes.